\documentclass[11pt,letterpaper]{article}
\usepackage[utf8]{inputenc}
\usepackage[T1]{fontenc}
\usepackage{lmodern}
\usepackage[letterpaper,top=2.2cm,bottom=2.2cm,left=2.0cm,right=2.0cm]{geometry}
\usepackage{amsmath}
\usepackage{booktabs}
\usepackage{tabularx}
\usepackage{xcolor}
\usepackage{fancyhdr}
\usepackage{tcolorbox}

\definecolor{navyblue}{RGB}{26, 79, 124}
\pagestyle{fancy}
\fancyhf{}
\lhead{\small Solvanta Engineering Solutions -- Matriz de Trazabilidad}
\rhead{\small Consistencia Inter-Documental DWSIM / TC-01 / TE-01 / PFD v0}
\cfoot{\small Página \thepage\ de 1}
\renewcommand{\headrulewidth}{0.4pt}

\begin{document}

\begin{center}
{\large \textbf{SOLVANTA ENGINEERING SOLUTIONS S.A. DE C.V.}}\\[0.2cm]
{\Large \textbf{\color{navyblue} MATRIZ DE CONSISTENCIA Y TRAZABILIDAD TÉCNICA}}\\[0.1cm]
{\small Expediente VTP-01 | Planta de e-Metanol ($30\,000\text{ t/año}$) | Revisión Oficial v1.0 Estandarizada}
\end{center}

\vspace{0.2cm}
\noindent La presente matriz garantiza la correspondencia unívoca y rigurosa entre las corrientes y equipos modelados en el simulador \textbf{DWSIM v1}, las tablas oficiales de materia (\textbf{TC-01}) y energía (\textbf{TE-01}), y la representación física de ingeniería en el diagrama \textbf{PFD v0}, estructurados por áreas funcionales ($100$ a $800$).

\vspace{0.2cm}
\begin{center}
\footnotesize
\begin{tabularx}{\textwidth}{lp{3.5cm}p{3.0cm}p{2.5cm}X}
\toprule
\textbf{Tag} & \textbf{Objeto en DWSIM v1} & \textbf{TC-01 / TE-01} & \textbf{Tag PFD v0} & \textbf{Resultado de Revisión Cruzada} \\
\midrule
\texttt{S-101} & Alimentación $\text{CO}_2$ captura & TC-01: \texttt{S-101} & Línea \texttt{S-101} & Consistente ($20\text{ bar}, 25^\circ\text{C}, 5\,222.4\text{ kg/h}$). \\
\texttt{EL-201}& Water Electrolyzer SOEC & TE-01: \texttt{E-SOEC} & \texttt{EL-201} & Consistente ($54.4\text{ MW}, 750^\circ\text{C}, 20\text{ bar}$). \\
\texttt{E-201} & Carcasa caldera acoplada & TE-01: \texttt{E-201} & Carcasa \texttt{R-401} & Consistente ($9\,139.2\text{ kW}$, vapor saturado $220^\circ\text{C}$). \\
\texttt{E-202} & Sobrecalentador vapor SOEC& TE-01: \texttt{E-202} & \texttt{E-202} & Consistente ($5\,119.2\text{ kW}$, $220 \to 750^\circ\text{C}$). \\
\texttt{C-301} & Compresor syngas multietapa & TE-01: \texttt{E-C301} & \texttt{C-301} & Consistente ($1\,138.7\text{ kW}, 19.4 \to 75.0\text{ bar}$). \\
\texttt{R-401} & Reactor multitubular BWR & TE-01: \texttt{E-R401} & \texttt{R-401} & Consistente (Lurgi, $245^\circ\text{C}, 75\text{ bar}$). \\
\texttt{E-502/503}& Economizador FEHE ($272\text{ kW}$) & TE-01: \texttt{E-502/E-503} & \texttt{E-502/E-503} & Consistente ($272.05\text{ kW}$, balance neto $0.00\text{ kW}$). \\
\texttt{E-501} & Condensador efluente síntesis& TE-01: \texttt{E-501} & \texttt{E-501} & Consistente ($13\,952.7\text{ kW}$, $240.4 \to 38.0^\circ\text{C}$). \\
\texttt{V-501} & Separador flash frío HP & TC-01: \texttt{S-501/S-507}& \texttt{V-501} & Consistente ($75\text{ bar}, 38^\circ\text{C}$, vapor/líquido). \\
\texttt{C-501} & Soplador gas de reciclo & TE-01: \texttt{E-C501} & \texttt{C-501} & Consistente ($150\text{ kW}$, gas seco a $43^\circ\text{C}$). \\
\texttt{X-501} & Turboexpansor purga continua & TE-01: \texttt{E-TURB} & \texttt{X-501} & Consistente ($-390.43\text{ kW}, 75 \to 1.1\text{ bar}$). \\
\texttt{V-601} & Flash desgasificador LP & TC-01: \texttt{S-601/S-602}& \texttt{V-601} & Consistente ($3\text{ bar}, 38^\circ\text{C}, 0.0\text{ kW}$). \\
\texttt{T-601} & Columna fraccionadora & TC-01: \texttt{S-604/S-605}& \texttt{T-601} & Consistente (28 platos, separación metanol/agua). \\
\texttt{S-604} & e-Metanol comercial Grado AA & TC-01: \texttt{S-604} & Línea \texttt{S-604} & Consistente ($3\,750.0\text{ kg/h}$ netos, $\ge 99.85\%$). \\
\texttt{V-801} & Flash desgasificador ZLD & TC-01: \texttt{S-801/S-802}& \texttt{V-801} & Consistente ($19.4\text{ bar}, 38.7^\circ\text{C}$, desgasificación). \\
\texttt{P-801} & Bomba retorno ZLD a caldera & TE-01: \texttt{E-P801} & \texttt{P-801} & Consistente ($0.20\text{ kW}, 19.4 \to 20\text{ bar}$). \\
\bottomrule
\end{tabularx}
\end{center}

\vspace{0.3cm}
\noindent\textbf{Protocolo de Control y Resolución de Discrepancias:}
\begin{enumerate}
    \item \textbf{Detección:} Si se identifica una divergencia numérica, se localiza el documento de origen.
    \item \textbf{Verificación:} Se confronta el valor contra el archivo congelado \texttt{DWSIM\_v1\_APROBADA\_Solvanta.dwxmz}.
    \item \textbf{Corrección:} Se actualiza la tabla técnica o el PFD correspondiente.
    \item \textbf{Documentación:} Se registra en la bitácora técnica de ingeniería.
    \item \textbf{Versionamiento:} Se incrementa el consecutivo de revisión del documento afectado.
\end{enumerate}

\vspace{0.5cm}
\begin{center}
\footnotesize
\begin{tabularx}{0.95\textwidth}{XcX}
\toprule
\textbf{Elaboró y Verificó:} & \textbf{Revisión Técnica:} & \textbf{Aprobación Docente:} \\
\midrule
Ing. Karol Paola Camacho López & Rev. v1.0 (Estandarizada) & Dr. Jesús Andrés Arzola Flores \\
Solvanta Engineering Solutions & 14 de septiembre de 2026 & DPyPQ-2026-1 \\
\bottomrule
\end{tabularx}
\end{center}

\end{document}